\SourceLeaf{71}
\[
\cos\varphi^r=(1+a)\frac{1}{2^r}.{}^{r}\mathfrak{B}^{(\frac12(r+a))}\cos a\varphi+.\,.\,.+\frac{1}{2^{r-1}}\mathfrak{B}^{(\frac12(r+a)+h)}\cos(2h+a)\varphi+.\,.\,.,
\]

\noindent und so lassen sich zwei Reihen, die auf den ersten Blick völlig verschieden schienen, in einem einzigen, recht einfachen Ausdruck zusammenfassen.

\Needspace{10\baselineskip}\par\smallskip\hrulefill\smallskip\par

\begin{center}\textbf{BEMERKUNGEN.}\end{center}

Das Manuskript der vorliegenden »Dissertatiuncula« wurde 1893 von Herrn Prof. \textsc{Wilhelm Meyer} (Göttingen) in den Akten der Königlichen Gesellschaft der Wissenschaften zu Göttingen gefunden und von E. \textsc{Schering} mit einer Reihe erläuternder Bemerkungen in den »Göttinger Nachrichten« von 1893, Nr.\,15, veröffentlicht. In der Urschrift fehlen der Name des Verfassers und eine Angabe zu Ort und Zeit der Abfassung. Der Vergleich anderer Manuskripte von \textsc{Gauss} mit der vorliegenden Handschrift beseitigte jedoch jeden Zweifel, der sonst vielleicht an \textsc{Gauss}' Autorschaft bestehen könnte.

Inhaltlich ist zu bemerken, dass \textsc{Gauss} vermutlich schon zu der Zeit, als er das Collegium Carolinum in Braunschweig besuchte (1792 bis 1795), die Werke von \textsc{Lagrange} und \textsc{Euler} studierte. Im Nachlass finden sich verschiedene Auszüge aus Werken von \textsc{Lagrange} und \textsc{Euler}, die wahrscheinlich aus dieser Zeit stammen, darunter insbesondere Auszüge aus \textsc{Euler}s Institutiones calculi integralis, Teil I, Kapitel 6, an die die Dissertatiuncula unmittelbar anknüpft.

Ist damit auch inhaltlich \textsc{Gauss}' Verfasserschaft unzweifelhaft, so lässt sich die fast vollständige Abhängigkeit der ganzen Entwicklung von \textsc{Euler} gleichwohl nur mit der Annahme vereinbaren, dass die fragliche Untersuchung ganz zu Beginn der Göttinger Studienzeit von \textsc{Gauss} (Herbst 1795 bis 1798) oder sogar noch davor entstand. \textsc{Euler} löst nämlich die Aufgabe, $(1+n\cos\varphi)^\nu$ in die Reihe

\[
A + A^{(1)}\cos\varphi + A^{(2)}\cos 2\varphi + A^{(3)}\cos 3\varphi + .\,.\,.
\]

\noindent zu entwickeln, indem er nur $A$ und $A^{(1)}$ direkt angibt und für die weiteren Koeffizienten auf eine Rekursionsformel verweist, die sich durch einen naheliegenden Kunstgriff gewinnen lässt. Ohne diesen Kunstgriff gibt \textsc{Gauss} auch die höheren Koeffizienten $A^{(2)}$,\,.\,. durch direkte Betrachtung an, übernimmt im Übrigen aber alle Ansätze und Weiterentwicklungen \textsc{Euler}s.

Die Verweise auf \textsc{Kästner} sprechen dafür, dass die endgültige Abfassung der Dissertatiuncula in den Beginn der Göttinger Studienzeit fällt. Die gegen Ende von Nr.\,10 genannte Abhandlung \textsc{Klügel}s gehört zu den mathematischen Kommentationen der Göttinger Sozietät von 1793 und 94. Diese wurden als Gesamtband allerdings erst 1796 herausgegeben, doch waren \textsc{Gauss} die einzelnen Abhandlungen wahrscheinlich schon früher zugänglich.

Für die Annahme, dass die Dissertatiuncula spätestens Anfang 1796 entstand, spricht auch, dass sich in dem Ende März 1796 beginnenden wissenschaftlichen Tagebuch keine Notiz über diese Abhandlung findet.

\Needspace{10\baselineskip}
Die Entwicklung unter Nr.\,8, $\gamma$, erfordert noch eine besondere Bemerkung. Der Vorwurf eines Fehlers bei \textsc{Euler} ist unberechtigt. Vielmehr hat sich \textsc{Gauss} versehen, indem er bei der Berechnung der Gleichungen für $A^{(h)}$, $A^{(h+1)}$ auf der rechten Seite den Faktor $\frac{2}{\BareRoot{(1-nn)}}$ übersah. Indem \textsc{Schering} in der von ihm besorgten Ausgabe diesen Faktor hinzufügte, stellte er die Übereinstimmung mit \textsc{Euler} her. Der vorliegende Abdruck gibt ebenfalls die berichtigten Formeln.\hfill\textsc{Fricke}.

\par\smallskip\hrulefill\smallskip\par

